\chapter{Volume, injectivity radius and injectivity angle}
\source{cheeger-gromov-taylor:page-0028,cheeger-gromov-taylor:page-0029,cheeger-gromov-taylor:page-0030,cheeger-gromov-taylor:page-0031,cheeger-gromov-taylor:page-0032,cheeger-gromov-taylor:page-0033,cheeger-gromov-taylor:page-0034,cheeger-gromov-taylor:page-0035,cheeger-gromov-taylor:page-0036,cheeger-gromov-taylor:page-0037,cheeger-gromov-taylor:page-0038,cheeger-gromov-taylor:page-0039,cheeger-gromov-taylor:page-0040}

Let $x\in M^n$, and let $A(r,\theta)$,
$\mathcal{Q}^H(r)$, $V_r^H$ be as defined after (1.7) and in (1.16), (1.17).
Suppose $\Ric_M\geq (n-1)H$ and regard $A(r,\theta)$ as a function on the tangent
space $T_xM$. If $H>0$, we restrict attention to $B_{\pi/\sqrt{H}}(0)\subset T_xM$.
Let $U\subset T_xM$ denote the interior of the cut locus in the tangent space.
The standard volume comparison theorem (see
\cite[pp.~253--257]{cgt-bishop-crittenden}) shows that
on $U$,
\begin{equation}
\tag{4.1}
\frac{A(r,\theta)}{\mathcal{Q}^H(r)}\downarrow,
\end{equation}
where $\downarrow$ denotes decrease. If $S_r(0)\subset T_xM$ denotes the sphere
of radius $r$ about the origin, (4.1) implies
\begin{equation}
\tag{4.2}
\frac{\int_{S_r(0)\cap U} A(r,\theta)}{\int_{S_r(0)} \mathcal{Q}^H(r)}\downarrow;
\end{equation}
Cut points can only decrease the numerator. More generally, for positive $f,g$ with
$f/g\downarrow$,
\begin{equation}
\tag{4.3}
\frac{\int_0^r f(s)}{\int_0^r g(s)}\downarrow.
\end{equation}
To see this, set $f/g=h$. Then $h\downarrow$ and
\begin{equation}
\tag{4.4}
\int_0^r f \int_r^R g
= \int_0^r gh \int_r^R g
\geq \left[\int_0^r g\right] h(r)\int_r^R g
\geq \left[\int_0^r g\right]\int_r^R gh
= \int_0^r g \int_r^R f.
\end{equation}
Thus, equivalently,
\begin{equation}
\tag{4.5}
\int_0^r f \int_0^R g
= \int_0^r f \int_0^r g + \int_0^r f \int_r^R g
\geq \int_0^r f \int_0^r g + \int_0^r g \int_r^R f
= \int_0^r g \int_0^R f,
\end{equation}
which gives (4.3). Applying (4.2) and (4.3) yields the following estimates.

\begin{proposition}[Relative volume estimates]
\label{prop:relative-volume-estimates}
\source{cheeger-gromov-taylor:page-0028,cheeger-gromov-taylor:page-0029}
\dcref{Proposition 4.1}
\group{section-4}

Let $M^n$ be complete and $\Ric_M\geq (n-1)H$. Let $p,x\in M^n$ and
$\rho(p,x)=\rho$. Then
\begin{equation}
\tag{i}
\frac{V_r(p)}{V_r^H}\downarrow,
\end{equation}
or equivalently, for $r_1 \le r_2 \le r_3$,
\begin{equation}
\tag{ii}
\frac{V_{r_3}(p)-V_{r_2}(p)}{V_{r_3}^H-V_{r_2}^H}\le \frac{V_{r_1}(p)}{V_{r_1}^H}.
\end{equation}
Moreover,
\begin{equation}
\tag{iii}
V_{r_2}(p)\frac{V_{r_1}^H}{V_{\rho+r_2}^H}\le V_{r_1}(x).
\end{equation}
If $r_2+r<\rho$, then
\begin{equation}
\tag{iv}
V_{r_2}(p)\frac{V_r^H}{V_{\rho+r_2}^H-V_{\rho-r_2}^H}\le V_r(x).
\end{equation}
\end{proposition}

\begin{proof}
\source{cheeger-gromov-taylor:page-0029}
\sketch
(i) and (ii). Let
\begin{equation}
\tag{4.10}
f=\int_{S_r^{n-1}(0)\cap U}A(r,\theta), \qquad g=\int_{S_r^{n-1}(0)}\mathcal{Q}^H(r).
\end{equation}
Then (i) and (ii) follow from (4.2), (4.3).

(iii) We have
\begin{equation}
\tag{4.11}
V_{r_2}(p)\le V_{\rho+r_2}(x)\le V_{\rho+r_2}^H\frac{V_{r_1}(x)}{V_{r_1}^H}.
\end{equation}

(iv) Using (ii) we obtain
\begin{equation}
\tag{4.12}
V_{r_2}(p)\le V_{\rho+r_2}(x)-V_{\rho-r_2}(x)\le \left(V_{\rho+r_2}^H-V_{\rho-r_2}^H\right)\frac{V_r(x)}{V_r^H}.
\end{equation}
\end{proof}

If $\Ric_M\ge 0$ and $d>2r$, Proposition~\ref{prop:relative-volume-estimates}(iv) gives
\begin{equation}
\tag{4.13}
\frac{(d-2r)^n}{d^n-(d-2r)^n}V_r(p)\le V_d(p).
\end{equation}

\begin{remark}[Relative volume comparison for geodesically star-shaped sets]
\label{rem:relative-volume-star-shaped}
\source{cheeger-gromov-taylor:page-0029}
\dcref{Remark 4.1}
\group{section-4}

Suppose $S\subset M^n$ contains the unique minimizing geodesic from $p$ to
each $q\in S$ outside the cut locus of $p$. Proposition~\ref{prop:relative-volume-estimates} remains valid with
$V_r(p)$ replaced by $V(S\cap B_r(p))$.
\end{remark}

\begin{remark}[Ricci models]
\label{rem:ricci-models}
\source{cheeger-gromov-taylor:page-0030}
\dcref{Remark 4.2}
\group{section-4}

For a Ricci model with metric $dr^2+f^2(r)g$, where
$f(r)\sim r-Kr^3/6+\cdots$, one has
\[
\Ric_{M^n}\geq -(n-1)f''(r)/f(r).
\]
Remark~\ref{rem:relative-volume-star-shaped} then holds with $V_r^H$ replaced by $f^{n-1}(r)$.
\end{remark}

We next consider the refined injectivity angle $\wtomega(H,\epsilon,r,x)$
under the assumptions of Proposition~\ref{prop:relative-volume-estimates}.

\begin{proposition}[Lower estimates for the refined injectivity angle]
\label{prop:refined-injectivity-angle-lower-estimates}
\source{cheeger-gromov-taylor:page-0030,cheeger-gromov-taylor:page-0031}
\dcref{Proposition 4.2}
\group{section-4}

\textup{(i)} Given $\epsilon$, $r_2$, $r$, we have
\begin{equation}\tag{4.14}
\wtomega(H,\epsilon,r,x)\ge \frac{1}{1-\epsilon}
\left[\frac{V_{r_2}(p)-V_r^H}{V^H_{\rho+r_2}-V_r^H}-\epsilon\right],
\end{equation}
where the right-hand side is $>0$, if $\epsilon,r$ are sufficiently small.

\textup{(ii)} If moreover $r_2<\rho$, then
\begin{equation}\tag{4.15}
\wtomega(H,\epsilon,\rho-r_2,x)\ge \frac{1}{1-\epsilon}
\left[\frac{V_{r_2}(p)}{V^H_{\rho+r_2}-V^H_{\rho-r_2}}-\epsilon\right].
\end{equation}

\textup{(iii)} Thus if $r(H,V)$ is defined by $\frac{V}{2}=V^H_{r(H,V)}$, then
\begin{equation}\tag{4.16}
\wtomega\!\left(H,\frac{V_{r_2}(p)}{2(2V^H_{\rho+r_2}-V_{r_2}(p))},r(H,V_{r_2}(p)),x\right)
\ge \frac{V_{r_2}(p)}{4V^H_{\rho+r_2}-3V_{r_2}(p)}.
\end{equation}

\textup{(iv)} If moreover $r_2<\rho$, then
\[
\wtomega\!\left(H,\frac{V_{r_2}(p)}{4(V^H_{\rho+r_2}-V^H_{\rho-r_2})},\rho-r_2,x\right)
\ge \frac{V_{r_2}(p)}{4(V^H_{\rho+r_2}-V_{\rho-r_2}^H)-V_{r_2}(p)}.
\]
\end{proposition}

\begin{proof}
\source{cheeger-gromov-taylor:page-0030,cheeger-gromov-taylor:page-0031}
\sketch
\textup{(i)} Since $V_r(x)\le V_r^H$ and $B_{r_2}(p)\subset B_{\rho+r_2}(x)$, we have
\begin{equation}\tag{4.17}
V_{r_2}(p)-V_r^H\le V_{\rho+r_2}(x)-V_r(x).
\end{equation}
Using (4.2) and integrating in polar coordinates, we easily get
\begin{equation}\tag{4.18}
V_{\rho+r_2}(x)-V_r(x)\le
\bigl(V^H_{\rho+r_2}-V_r^H\bigr)\bigl[\wtomega(H,\epsilon,r,x)+\epsilon(1-\wtomega(H,\epsilon,r,x))\bigr],
\end{equation}
from which (i) follows.

\textup{(ii)} This follows as in (i), by noting that
\begin{equation}
\tag{4.19}
V_{r_2}(p) \leq V_{\rho+r_2}(x) - V_{\rho-r_2}(x).
\end{equation}

\textup{(iii)} and (iv) follow immediately from (i) and (ii).
\end{proof}

For a shortest geodesic loop $\gamma$ at $p$ with $L[\gamma]=2l$ and
$K_M\leq K$ on $B_l(p)$, Klingenberg's estimate gives
$i(p)\geq\min(l,\pi/\sqrt K)$, with the second term interpreted as $\infty$
when $K\leq0$ \cite{cgt-cheeger-ebin}.

For a compact manifold with $d(M^n)\leq d$, $V(M^n)\geq V$, and
$K_M\geq H$, the finiteness theorem \cite{cgt-cheeger-finiteness} gives a
constant $c_n(d,V,H)>0$ and
\begin{equation}
\tag{4.20}
i(p)\geq \min\left(\frac12c_n(d,V,H),\pi/\sqrt K\right).
\end{equation}

The local result below replaces this global compactness hypothesis by a lower
bound for the volume of one ball, while retaining a lower bound for the
distance to the conjugate locus.

To emphasize the local nature of our result, we now let $B_r(p)$ be a metric
ball in a Riemannian manifold such that for $r' < r$, $\overline{B_{r'}(p)}$ is
compact. Assume that on $B_r(p)$, $K_M \le K$ and $r \le \pi/\sqrt{K}$, ($r$
arbitrary if $K \le 0$). We let $V_r^0(p)$ denote the volume of $B_r(0) \subset
T_pM$ with respect to the pulled back metric.

\begin{theorem}[Local injectivity-radius bound]
\label{thm:local-injectivity-radius-bound}
\source{cheeger-gromov-taylor:page-0032,cheeger-gromov-taylor:page-0033}
\dcref{Theorem 4.3}
\group{section-4}

Let $B_r(p)$ be as above. If $\gamma$ is a geodesic loop on $p$ of length
$2l$, and $r_0+2s \le r$, $r_0 \le r/4$, then
\begin{equation}\tag{4.21}
l \ge \frac{r_0}{2}\,\frac{1}{1+V^0_{r_0+s}(p)/V_s(p)}.
\end{equation}
If in addition $H < K_M \le K$, then
\begin{equation}\tag{4.22}
l \ge \frac{r_0}{2}\,\frac{1}{1+V^H_{r_0+s}(p)/V_s(p)}.
\end{equation}
\end{theorem}

\begin{proof}
\source{cheeger-gromov-taylor:page-0033}
\sketch For a loop with $L(\gamma)=2l<r_0$, set
$N=\lfloor r_0/(2l)\rfloor$. Lemmas~\ref{lem:radial-class-concatenation},
\ref{lem:even-covering}, and \ref{lem:iterated-geodesic-loops} give at least $N$ inverse images
of each point of $B_s(p)$ in $B_{r_0+s}(0)$. Since $\exp_p$ is nonsingular and
orientation preserving, $N V_s(p)\leq V^0_{r_0+s}(p)$; rearranging and using
volume comparison yields (4.21) and (4.22).
\end{proof}

For $r$ below the conjugate radius, the Gauss lemma identifies $B_r(0)\subset
T_pM$ with radial geodesics from $p$. Write $c_1\sim_Ac_2$ when the curves
are homotopic relative endpoints through curves of length at most $A$; the
radial lift of a curve of length $<r$ is unique. Since $\exp_p$ is nonsingular
on $B_r(0)$, distinct radial geodesics represent distinct such classes. Write
$[c]$ for the radial representative and $[p]$ for the constant loop.

Let $\theta_1,\theta_2$ be closed curves on $p$, and let $c$ be an arc from $p$
with $L[\theta_i]\le 2l$, $L[c]=s$ and $2(l+s)<r$. Let $\theta_i \cup c$ denote
the arc from $p$ which is equal to $\theta_i$ followed by $c$.

\begin{lemma}[Concatenation preserves radial classes]
\label{lem:radial-class-concatenation}
\source{cheeger-gromov-taylor:page-0032}
\dcref{Lemma 4.4}
\group{section-4}

If $2(l+s)<r$ and $[\theta_1\cup c]=[\theta_2\cup c]$, then $[\theta_1]=[\theta_2]$.
\end{lemma}

\begin{proof}
\source{cheeger-gromov-taylor:page-0032}
\sketch
If $[\theta_1\cup c]=[\theta_2\cup c]$, then $\theta_1\cup c \sim_{2(l+s)}
\theta_2\cup c$. Thus
\[
\theta_1 \sim_{2(l+s)} \theta_1\cup c\cup -c \sim_{2(l+s)}
\theta_2\cup c\cup -c \sim_{2(l+s)} \theta_2.
\]
By uniqueness of radial representatives, $[\theta_1]=[\theta_2]$.
\end{proof}

Let $\#(q,a)$ denote the number of distinct inverse images of $q \in B_r(p)$ in
$B_a(0)$, $a < r$.

\begin{lemma}[Even covering lemma]
\label{lem:even-covering}
\source{cheeger-gromov-taylor:page-0033}
\dcref{Lemma 4.5}
\group{section-4}

If $r_0 + 2s < r/4$ and $q \in B_s(p)$, then
\[
\#(q,3r_0) \ge \#(p,r_0).
\]
\end{lemma}

\begin{proof}
\source{cheeger-gromov-taylor:page-0033}
\sketch
By radial lifting, $\#(p,r_0)=m$, where $\gamma_1 \cdots \gamma_m$ are the distinct
geodesic loops on $p$ of length $< r_0$. Let $\sigma$ be a minimal geodesic from
$p$ to $q$. By Lemma~\ref{lem:radial-class-concatenation}, $[\gamma_i \cup \sigma]$, $i=1,\cdots,m$, represent
distinct inverse images of $q$ in $B_{r_0+s}(p)$.
\end{proof}

Let $\theta$ be a loop on $p$, and $\sigma$ a geodesic arc from $p$ with
$L(\theta)=2l$, $L(\sigma)=s$ and $2(l+s)<r$. The correspondence
$\sigma \mapsto [\theta \cup \sigma]$ defines a map $\pi_{[\theta]}: B_s(0) \to
B_{2l+s}(0)$ which clearly depends only on $[\theta]$. It is also obvious that
$\exp_p \pi_{[\theta]}(x)=\exp_p x$ for $x \in B_s(0)$. Hence $\pi_{[\theta]}$ is
an isometry from $B_s(0)$ to $\pi_{[\theta]}(B_s(0))$ with respect to the pulled
back metric on $B_{2l+s}(0)$. It follows immediately from Lemma~\ref{lem:radial-class-concatenation} that
$\pi_{[\theta]}$ has no fixed points unless $[\theta]=p$.

Set $i\gamma=\gamma\cup \cdots \cup \gamma$ and $L(\gamma)=2l$.

\begin{lemma}[Iterated geodesic loops are distinct]
\label{lem:iterated-geodesic-loops}
\source{cheeger-gromov-taylor:page-0033}
\dcref{Lemma 4.6}
\group{section-4}

Let $\gamma$ be a geodesic loop on $p$ with $2l \cdot N < \frac14 r$. Then
$[\gamma], [2\gamma], \cdots, [N\gamma]$ are all distinct.
\end{lemma}

\begin{proof}
\source{cheeger-gromov-taylor:page-0033}
\sketch
Since $(i+j)\gamma=i\gamma \cup j\gamma$, if
$[(i+j)\gamma]=[i\gamma \cup j\gamma]=[j\gamma]=[j\gamma \cup p]$, then
$[i\gamma]=p$ by Lemma~\ref{lem:radial-class-concatenation}. Regard $[p],[\gamma],\cdots,[(i-1)\gamma]$ as points
in the strictly convex ball $B_{2il}(0)$. Then $\pi_{[\gamma]}: B_{2il}(0) \to
B_{2(i+1)l}(0)$ preserving distance. Hence $\pi_{[\gamma]}$ preserves the unique
center of gravity of the set $\{[p],[\gamma],\cdots,[(i-1)\gamma]\}$. By
definition the center of gravity is the unique minimum of the function
$y \mapsto \sum_j \rho^2([j\gamma],y)$. It exists because $2il < \frac12
\pi/\sqrt{K}$ implies that $B_{2il}(0)$ is convex. This contradicts the fact that
$\pi_\gamma$ has no fixed points.
\end{proof}

If we combine Theorem~\ref{thm:local-injectivity-radius-bound} with Proposition~\ref{prop:relative-volume-estimates}, we immediately obtain the
following lower bound for the injectivity radius.

\begin{theorem}[Injectivity radius from relative volume]
\label{thm:injectivity-radius-relative-volume}
\source{cheeger-gromov-taylor:page-0033,cheeger-gromov-taylor:page-0034}
\dcref{Theorem 4.7}
\group{section-4}

Let $M^n$ be complete with $H \le K_M \le K$.
Let $\rho=\rho(p,x)$, and fix $r, r_0, s$, with $r_0 + 2s < \pi/\sqrt{k}$,
$r_0 \le \pi/4\sqrt{k}$.
\[
\text{(i)}
\]
\begin{equation}
\tag{4.23}
i(x) \ge \frac{r_0}{2}\,\frac{1}{1+\left(V^H_{r_0+s}/V_r(p)\right)\left(V^H_{\rho+r}/V^H_s\right)}.
\end{equation}
\[
\text{(ii)}
\]
Moreover, if $r+s<\rho$, then
\begin{equation}
\tag{4.24}
i(x)\ge \frac{r_0}{2}\,\frac{1}{1+\bigl(V^{H}_{r_0+s}/V_r(p)\bigr)\bigl(V^{H}_{\rho+r}-V^{H}_{\rho-r}\bigr)/V^{H}_s}.
\end{equation}
\end{theorem}

As an application of the relative volume estimates proved at the beginning of
this section, we now consider complete manifolds whose Ricci curvature at
infinity is ``almost nonnegative'' in the following sense. Fix \(p\in M^n\) and
set \(r=\rho(p,x)\). We assume that for some \(r_0>0\) and all \(r\ge r_0\),
\begin{equation}\tag{4.25}
\operatorname{Ric}_M(x)\ge (n-1)\frac{\frac14\pm v^2}{r^2}.
\end{equation}

(For Theorem~\ref{thm:asymptotic-ricci-compactness} and Theorem~\ref{thm:asymptotic-ricci-volume-growth}(iii) it suffices to assume that (4.25) holds
only for radial directions from \(p\); however for Theorem~\ref{thm:asymptotic-ricci-volume-growth}(i) and (ii) this
does not suffice.) By convention we take $v\ge 0$. Observe that the Jacobi
equation
\begin{equation}\tag{4.26}
g''=-\frac{\bigl(\frac14-v^2\bigr)g}{r^2}
\end{equation}
corresponding to sectional curvature \(\bigl(\frac14-v^2\bigr)r^2\) admits the
pair of solutions
\begin{equation}\tag{4.27}
g_{\pm v}(r)=r^{\frac12\pm v},
\end{equation}
and the particular solution
\begin{equation}\tag{4.28}
J_{v,R}=\frac{1}{2v}\Bigl(-r^{\frac12+v}+R^{2v}r^{\frac12-v}\Bigr),
\end{equation}
satisfying \(J_{v,R}(R)=0,\ J'_{v,R}(R)=-1\). For the case
\(\frac14+v^2\), we replace $v$ by $iv$ in (4.27), and obtain the solutions
\begin{equation}\tag{4.29}
r^{\frac12}\cos(v\log r),\qquad r^{\frac12}\sin(v\log r).
\end{equation}
Corresponding to (4.28) we have
\begin{equation}\tag{4.30}
\mathcal G_{iv,R}(r)=\frac{R}{v}\,r^{\frac12}\sin(v\log r/R).
\end{equation}
Note that
\begin{equation}\tag{4.31}
\mathcal G_{iv,R}(\Re^{\pi/v})=\mathcal G_{iv,R}(R)=0.
\end{equation}
This gives the following extension of Myer's theorem.

\begin{theorem}[Asymptotic Ricci lower bound and compactness]
\label{thm:asymptotic-ricci-compactness}
\source{cheeger-gromov-taylor:page-0034,cheeger-gromov-taylor:page-0035}
\dcref{Theorem 4.8}
\group{section-4}

Let \(M^n\) be a complete Riemannian manifold such that for some
$p\in M^n,\, r_0>0,$ and $v>0$ we have
\begin{equation}\tag{4.32}
\operatorname{Ric}_M(x)\ge (n-1)\frac{\frac14+v^2}{r^2}
\end{equation}
for all \(r\ge r_0\). Then \(M^n\) is compact and the diameter $d_p$ from $p$
satisfies
\begin{equation}\tag{4.33}
d_p<e^{\pi/v}r_0.
\end{equation}
\end{theorem}

\begin{proof}
\source{cheeger-gromov-taylor:page-0035}
\sketch
We use the field $\mathcal{F}_{\gamma|_{[0,e^{\nu}/r_0]}},$ and argue by index
comparison as in the standard proof of Myers' theorem. If
$\gamma|_{[0,e^{\nu}/r_0]}$ is minimal, (4.31) and (4.32) lead to the conclusion
that $\gamma|_{[r_0,e^{\nu}/r_0]}$ contains a pair of conjugate points.
\end{proof}

We now consider the case $(\frac14-\nu^2)/r^2$. Observe that for $\nu>0$, the
conditions
\begin{equation}
\nu \le \frac{n+1}{2(n-1)}, \qquad
-\frac{n}{(n-1)^2} \le \frac14-\nu^2, \qquad
0 \le (n-1)\left(\frac12-\nu\right)+1
\tag{4.36}
\end{equation}
are equivalent to one another. Also recall that if $M^n$ is not compact, a
normal geodesic $\gamma:[0,\infty)\to M^n$ is called a ray if each segment of
$\gamma$ is minimal. Parts (i) and (ii) of the following theorem give a lower
bound on the rate of growth of the volume of a ball, which generalizes the
previously mentioned result of Yau and Calabi for the case $\Ric_M\ge 0$.
A novel feature of the argument, as distinct from the case
$\Ric_M\ge (n-1)H$, will be the use of a family of model spaces with metric
\[
dr^2 + J^2_{\nu,R}(r)g,\qquad 0\le r\le R.
\]
Here the metric is expressed in polar coordinates, but the origin is at $R$.
For $\nu<\frac12$, the metric is singular at the boundary $r=0$, while for
$\nu>\frac12$ there is a conjugate point at the boundary. For $\nu=\frac12$
we have the usual flat metric on a ball of radius $R$.

\begin{theorem}[Asymptotic Ricci lower bounds and volume growth]
\label{thm:asymptotic-ricci-volume-growth}
\source{cheeger-gromov-taylor:page-0035,cheeger-gromov-taylor:page-0036,cheeger-gromov-taylor:page-0037}
\dcref{Theorem 4.9}
\group{section-4}

Let $M^n$ be a complete noncompact Riemannian manifold, and let
$\gamma:[0,\infty)\to M^n$ be a ray from some $p\in M^n$.
\begin{enumerate}
\item[(i)] Fix $0\le \nu\le \tfrac12(n+1)/(n-1)$. If $\nu\ge \tfrac12$, assume
that for some $r_0>0$ and all $r>r_0$ we have
\begin{equation}
\tag{4.37}
\Ric_M(x)\ge (n-1)\frac{(\frac14-\nu^2)}{r^2}.
\end{equation}

If $0 \leq \nu < \frac12$, assume that for some $r_1$ such that $r_0 < r_1 <
\frac32 r_0$ and for all $r > r_0$,
\begin{equation}\tag{4.38}
\RicM(x) \geq (n-1)\frac{\left(\frac14-\nu^2\right)}{\bigl(r-2(r_1-r_0)\bigr)^2}.
\end{equation}

If $0 \leq \nu < \frac12 (n+1)/(n-1)$, then there exists $C_{r_1}$,
independent of $M^n$ such that for $r > 2r_1-r_0$
\begin{equation}\tag{4.39}
c_{r_1}V_{(r_1-r_0)}(\gamma(r_1))\,r^{(n-1)(\frac12-\nu)+1} \leq V_{r+(r_1-r_0)}(p).
\end{equation}

\item[(ii)] If $\nu = \frac12 (n+1)/(n-1)$, and (4.37) holds, then for
$r > 2r_1-r_0$
\begin{equation}\tag{4.40}
c_{r_1}V_{(r_1-r_0)}(\gamma(r_1))\log r \leq V_{r+(r_1-r_0)}(p).
\end{equation}

\item[(iii)] For all $\nu \geq 0$, if (4.37) holds (in radial directions), then
\begin{equation}\tag{4.41}
V_r(p) < c_M r^{(n-1)(\frac12+\nu)+1}.
\end{equation}
\end{enumerate}
\end{theorem}

\begin{proof}
\source{cheeger-gromov-taylor:page-0036,cheeger-gromov-taylor:page-0037}
\uses{prop:relative-volume-estimates}
\sketch Apply Proposition~\ref{prop:relative-volume-estimates} to the model
Jacobi field determined by $dr^2+J_{\nu,R}^2(r)g$. For $\nu\geq 1/2$, the
curvature hypothesis gives the required radial comparison directly, and the
asymptotics of $J_{\nu,R}$ give the power lower bound (4.39). For
$0\leq\nu<1/2$, the shifted denominator in (4.38) gives the same comparison
after moving the model origin by $2(r_1-r_0)$. The endpoint value
$\nu=\frac12(n+1)/(n-1)$ is the critical exponent, hence the integral gives the
logarithmic lower bound (4.40). The upper bound (4.41) is the corresponding
Heintze--Karcher radial comparison \cite{cgt-heintze-karcher}.
\end{proof}

\begin{remark}[Sharpness of the volume-growth estimates]
\label{rem:sharpness-volume-growth-estimates}
\source{cheeger-gromov-taylor:page-0037}
\dcref{Remark 4.4}
\group{section-4}

The metrics on $\R^n$ that outside a compact set have the form
\begin{equation}
dr^2+r^{1\pm 2v}g
\tag{4.51}
\end{equation}
show that Theorem~\ref{thm:asymptotic-ricci-volume-growth} is sharp. Their injectivity angle decays no faster than
$r^{-(n-1)(\frac12+v)}$.
\end{remark}
